\section*{Abstract}

Diese Arbeit behandelt einen Sicherheitspatch für das Open-Source-Projekt
dnchub-app (FleetOptima), ein mandantenfähiges Flottenmanagement-System.
Zwei Issues im Repository des Projekts liessen sich auf dieselbe Stelle
im Backend zurückführen, nämlich die generische Zugriffsklasse
\texttt{BaseService} innerhalb der Python- und FastAPI-Codebase. Weil dort
beim Einzelzugriff per UUID weder die \texttt{organiza\-tion\_id} noch
das \texttt{deleted\_at}-Flag berücksichtigt wurden, konnte jeder
angemeldete Benutzer Datensätze fremder Mandanten lesen, verändern und
löschen, wie in Issue \#5 beschrieben \parencite{github_issue_5}.
Ausserdem konnten soft-gelöschte Einträge weiterhin per UUID
aufgerufen werden, wie Issue \#14 belegt \parencite{github_issue_14}.

Der Patch setzt direkt an der Wurzel an, indem die Basisklasse zwei
zusätzliche Parameter erhält und die fehlenden Filter dort einmalig
einzieht. Dadurch wird mit einer kleinen Änderung über 90 abhängige
Endpunkte abgesichert. Zusätzlich entstehen automatisierte Integrationstests,
die den Angriff gegen den ursprünglichen Code nachstellen und nach dem
Fix als Regressionsschutz dienen.

\vspace{2ex}

\noindent Keywords: IDOR, CWE-639, Multi-Tenant Authorization,
Soft-Delete Bypass, FastAPI, SQLAlchemy, Security Patch

\clearpage

\section*{Vorwort}

Dieses Dokument entstand im Modul Cyber Sicherheit (CYSL) an der FHNW
Hochschule für Technik im Frühlingssemester 2026. Aufgabe war es, für ein
reales Open-Source-Projekt auf GitHub einen
Sicherheitspatch einzureichen. Nach Abklärung mit dem Dozenten fiel
die Wahl auf dnchub-app. Das Projekt brachte mehrere offene und
konkret beschriebene Security-Issues mit, eine realistisch gewachsene
Codebase und einen aktiven Repository-Owner. Das Repository wurde
geforkt, sämtliche Änderungen liegen auf einem eigens dafür angelegten
Feature-Branch.

